\subsection{A direct genus-zero stabilizer and transporter theorem}

\begin{result}{E.41}{Proved}{Genus-zero retractions, normalizers and all finite transporter quotients}
Let $G=\Gamma_{0,n}$ be the pure labelled mapping class group,
$n\geq3$, and let $\sigma$ be an actual multicurve.  Write
$C_\sigma=\Stab_G(\sigma)$ and let $\Delta_\sigma$ be its
positive twist lattice.  There is an explicit retraction
$G\to C_\sigma$.  In particular $C_\sigma$ is closed in the
profinite topology of $G$, has its full induced topology, and
\[
 N_{\widehat G}(\widehat\Delta_\sigma)
 =C_{\widehat G}(\widehat\Delta_\sigma)
 =\widehat C_\sigma.                                         \tag{E.170}
\]
For the entire genus-zero exit category, the finite ambient
transporter quotients $C_N$ below, with $N$ ranging over all
finite-index normal subgroups of $G$, are cofinal among all
finite categorical congruences.  Composition with a fixed
crossing from a target stratum to the open stratum consequently
embeds every completed Hom-set into
$\widehat G/\widehat\Delta_\sigma$.
These statements prove the full exit-category comparison used
in E.16, for every $n$, independently of a procongruence curve
or multitwist normalizer theorem.
\end{result}
\begin{proof}
The stable dual graph is a tree with its $n$ leaves individually
labelled.  An automorphism of such a stable tree fixing every
leaf is the identity: each edge is determined by the two leaf
sets on its sides, and each vertex by its incident edges and
legs.  The two sides of any edge are also distinguished by
their nonempty, disjoint leaf sets.  Thus $C_\sigma$ fixes every
curve, side and cut component.  Cutting and capping its boundary
circles by marked discs gives the exact sequence
\[
 1\longrightarrow\mathbb Z^{E(\sigma)}
 \longrightarrow C_\sigma
 \xrightarrow{\ f\ }
       \prod_{v}\Gamma_{0,\operatorname{val}(v)}
 \longrightarrow1.                                         \tag{E.171}
\]
Its kernel consists precisely of the collar twists, and is
central.  Indeed choose boundary-fixed representatives of the
mapping classes on every cut component and glue them; changing
such representatives changes only the matched boundary twists.
Conversely a map trivial on every capped component is isotopic
to a product of those collar twists.  This proves the sequence
with its actual restriction and extension maps.

For each edge $e=(v,w)$, choose two leaves in two distinct
non-$e$ branches or legs at $v$, and two in two distinct
non-$e$ branches or legs at $w$.  Stability supplies the two
choices at each vertex.  Forgetting every other label gives a
four-point sphere on which $e$ survives.  No other edge of
$\sigma$ survives: in the minimal subtree joining these four
leaves the two branching vertices are exactly $v,w$, and the
path between them is the single edge $e$.  Every other edge
separates at most one selected leaf from the remaining three.
The resulting homomorphism to $\Gamma_{0,4}=F_2$ therefore
sends $\tau_e$ to its positive primitive cusp generator and
sends every $\tau_d$, $d\ne e$, to $1$.  Choose an integral
character of $F_2$ taking that cusp generator to $1$.  Its
composite gives a homomorphism $h_e:G\to\mathbb Z$ with
$h_e(\tau_d)=\delta_{ed}$.

Each component projection in (E.171) also extends to all of
$G$.  At the vertex $v$, choose one leaf in each branch incident
to $v$, retaining the given label for a leg at $v$.  Forgetting
the other labels collapses every outward branch to the chosen
mark and leaves precisely the component at $v$ with one mark
for each flag.  On $C_\sigma$ the resulting homomorphism
$f_v:G\to\Gamma_{0,\operatorname{val}(v)}$ is the component
projection of (E.171).  Use this actual collapse to choose its
basepoint identification; its restriction is then the stated
projection, including the node labels.  Put $h=(h_e)_e$ and
$f=(f_v)_v$.  We have constructed the global homomorphism
\[
 (h,f):G\longrightarrow
  \mathbb Z^{E(\sigma)}\times
                      \prod_v\Gamma_{0,\operatorname{val}(v)}.
                                                               \tag{E.172}
\]
Its restriction to $C_\sigma$ is an isomorphism.  It is
injective by (E.171) and the identity matrix on the twist
kernel.  For surjectivity, lift a prescribed component tuple
using (E.171), then multiply by the unique twist tuple which
corrects its $h$ coordinate.  Hence
\[
 r_\sigma=((h,f)|_{C_\sigma})^{-1}\circ(h,f):G\to C_\sigma
                                                               \tag{E.173}
\]
is the required retraction.

Any finite quotient of $C_\sigma$ extends to $G$ by composition
with $r_\sigma$, proving the full induced topology.  The group
$G$ is residually finite, as follows for example from its pure
sphere configuration free-group towers.  For $x\notin
C_\sigma$, choose a finite quotient $a:G\to F$ separating
$x$ from $r_\sigma(x)$.  The homomorphism
$(a,a r_\sigma):G\to F\times F$ sends $C_\sigma$ into the
diagonal and sends $x$ outside it.  This proves separability.
It also proves that $\widehat C_\sigma$ embeds as the closure
of $C_\sigma$ in $\widehat G$.  The characters $h_e$ prove
that the full completed twist lattice embeds there as well.

We give the independent fixed-subgroup input needed for the
normalizer assertion.  Let $X,Y$ be finitely generated free
groups, let $w$ be primitive in $X$, and let $T$ be the
automorphism of $F=X*Y$ with
$T(x)=x$ for $x\in X$ and $T(y)=wyw^{-1}$ for $y\in Y$.
Then
\[
                 \operatorname{Fix}_{\widehat F}(T)
                         =\widehat X.                         \tag{E.174}
\]
To prove this, write $X=\langle w\rangle*X_0$ and form
$E=F\rtimes_T\langle t\rangle$.  Put $u=w^{-1}t$.
Eliminating $w=t u^{-1}$ gives the exact presentation
\[
 E=\langle t,u,X_0,Y\mid [t,u]=1,\ [t,x]=1\ (x\in X_0),
                                  \ [u,y]=1\ (y\in Y)\rangle.
                                                               \tag{E.175}
\]
Here $X_0,Y$ denote free bases.  This is a right-angled Artin
group.  Its discrete centralizer of $t$ is
$A=X\times\langle t\rangle$.  For an elementary check, write
$E=A*_{\langle u\rangle}(Y\times\langle u\rangle)$.
The element $t$ fixes the $A$ vertex of its ordinary
Bass--Serre tree and no adjacent edge.  Otherwise, since $t$
is central in $A$, one would have $t=u^k$ in $A$; the
independent $w,t$ coordinates in $A^{\mathrm{ab}}$ exclude
that equality.  If $t$ fixed a different vertex it would fix
the path and its first edge, again impossible.  Its unique
fixed vertex therefore forces its centralizer into $A$, and
all of $A$ does centralize it.

Minasyan's Theorem 1.1 proves hereditary conjugacy separability
for right-angled Artin groups; Proposition 3.2 gives its
centralizer condition, and Proposition 12.1 identifies that
condition with
$C_{\widehat E}(t)=\overline{C_E(t)}$
\cite[Theorem 1.1; Proposition 3.2; Proposition 12.1,
pp.~383--384]{ET_Minasyan2012}.
These exact results and proofs apply to the finite graph in
(E.175).  The subgroup $A$ is a retract, by killing the $Y$
generators in that presentation, so its closure is its full
completion.  The subgroup $F$ also has full induced topology
in $E$: given a finite-index subgroup of $F$, choose a smaller
characteristic finite-index subgroup $L$; the automorphism of
$F/L$ induced by $T$ has finite order $m$, and the quotient
$(F/L)\rtimes C_m$ of $E$ realizes this finite quotient on $F$.
Consequently $\widehat F$ embeds in $\widehat E$, and right
exactness of completion identifies it with the kernel of
$\widehat E\to\widehat{\mathbb Z}$, $t\mapsto1$.  Intersecting
$C_{\widehat E}(t)=\widehat X\times\widehat{\mathbb Z}$ with
that kernel proves (E.174).

We next prove the singleton case of (E.170) by induction on
$n$.  A curve enclosing two labels is the rational pair-tail
case of E.33; its proof uses the full Birman sequence and the
peripheral self-normalizer proved there.  This includes every
curve when $n=4$.  If both sides of $e$ contain at least three
labels, forget a label $p$ on one side.  The image $e'$ remains
essential.  A normalizer element $q\in\widehat G$ sends
$\tau_e$ to $\tau_e^\lambda$ for a unit
$\lambda\in\widehat{\mathbb Z}^{\times}$.  Its image in the
smaller completed group belongs to $\widehat C_{e'}$ by
induction, and so centralizes $\tau_{e'}$.  Since the latter
has full infinite procyclic topology, $\lambda=1$.
The discrete map $C_e\to C_{e'}$ is onto, by forgetting $p$
on its cut component in (E.171), and its completed map is onto.
Multiplying $q$ by an appropriate element of $\widehat C_e$
therefore leaves an element of the completed point-pushing
kernel $\widehat\pi_1(S_{0,n-1},p)$ which centralizes $\tau_e$.

Take the basepoint on the chosen side of $e$.  A standard
peripheral basis gives
$\pi_1(S_{0,n-1},p)=X*Y$, where $X$ is the fundamental group
of that side of the cut surface, and $Y$ contains all but one
of the peripheral generators on the other side.  The boundary
word $w$ is the product of the distinct generators of $X$ and
is primitive: replacing its last generator by their product
is a free-basis change.  The twist $\tau_e$ acts on point
pushing by the partial conjugation $T$ in (E.174), after
choosing the positive collar direction (its inverse has the
same fixed subgroup).  This follows by carrying each based
peripheral loop across the collar: the inside loops do not
cross it; the outside loops acquire the entering $w$ and
returning $w^{-1}$.  Equation (E.174) makes the remaining
normalizer element lie in $\widehat X$, the completion of the
point pushes supported on the cut side.  That subgroup is
contained in $\widehat C_e$.  This proves
$N_{\widehat G}(\widehat{\langle\tau_e\rangle})
=\widehat C_e$ and also its equality with the centralizer.
The full completed Birman sequence used here is the actual
sequence established in E.33.

For a multicurve, an element normalizing
$\widehat\Delta_\sigma\cong\widehat{\mathbb Z}^{E(\sigma)}$
induces an invertible matrix on its twist basis.  Apply each
four-leaf forgetting map used above.  Its image normalizes the
primitive cusp subgroup of $\widehat\Gamma_{0,4}$, hence
centralizes that subgroup by the free-group peripheral
normalizer calculation of E.33.  Since this forgetting map
kills every other twist, its conjugation matrix row is exactly
the corresponding identity row.  The matrix is therefore the
identity, and the normalizer element centralizes every twist.
The singleton result puts it in every $\widehat C_e$.

For completeness the intersection of these subgroups is exactly
$\widehat C_\sigma$; this too follows without a profinite curve
complex theorem.  Induct on the number of curves, cutting first
along one $e$.  By (E.172) its stabilizer has the full completed
product description
$\widehat C_e=\widehat{\mathbb Z}\times
\widehat\Gamma_{0,A\cup\{*\}}\times
\widehat\Gamma_{0,B\cup\{*\}}$.
For another curve $d$ lying in the $A$ component, forgetting
all $B$ labels but its chosen node representative sends $C_d$
into the stabilizer of the surviving $d$ in that $A$ component.
Therefore an element of every $\widehat C_d$ has its $A$
projection in every corresponding completed stabilizer.
Induction identifies their intersection with the completed
multicurve stabilizer inside that component.  The same holds
on the $B$ side.  The preimage of these two component
stabilizers in $C_e$ is exactly $C_\sigma$; under the product
splitting it is the product with the unrestricted central
$\mathbb Z$.  The retractions already proved give the full
completion of both component stabilizers, so the same equality
holds after completion.  This proves the asserted intersection
and finishes (E.170).

We finally prove the full categorical cofinality.  Fix a finite
orbit skeleton with actual representatives $\sigma_i$.  For a
finite-index normal subgroup $N\lhd G$, let $q_N:G\to G/N$
and define
\[
 \Hom_{C_N}(i,j)=
   q_N\{g:g\sigma_i\supseteq\sigma_j\}/q_N\Delta_{\sigma_i}.
                                                               \tag{E.176}
\]
Composition is multiplication in $G/N$.  It is well-defined:
for an actual transporter $g:i\to j$ one has
$g^{-1}\Delta_{\sigma_j}g\subseteq\Delta_{\sigma_i}$, so
changing the left arrow by a source twist changes the product
only by a twist at its source.  All identities and associativity
are inherited from multiplication.  Thus $C_N$ is an actual
finite categorical quotient with the fixed object set.

For each pair $i,j$, list the finitely many submulticurves
$\rho\subseteq\sigma_i$ in the orbit of $\sigma_j$, and choose
$g_\rho$ with $g_\rho\rho=\sigma_j$.  We have the disjoint
transporter decomposition
\[
 \{g:g\sigma_i\supseteq\sigma_j\}
                 =\coprod_\rho C_{\sigma_j}g_\rho.            \tag{E.177}
\]
Indeed the subset $g^{-1}\sigma_j$ is exactly $\rho$, and two
transports of that subset differ on the left by
$C_{\sigma_j}$.  Since every genus-zero pure stabilizer fixes
its individual curves, right multiplication by
$\Delta_{\sigma_i}$ preserves each displayed coset.  In
particular $g_\rho\Delta_{\sigma_i}g_\rho^{-1}
\subseteq C_{\sigma_j}$.

Let $R$ be any finite categorical congruence.  The action of
$C_{\sigma_j}$ on its finitely many incoming Hom-sets modulo
$R$ has a finite-index normal kernel $M_j$.  The retraction
$r_{\sigma_j}$ gives a finite-index normal subgroup
$r_{\sigma_j}^{-1}(M_j)$ of $G$ whose intersection with
$C_{\sigma_j}$ is $M_j$.  Also, separability of
$C_{\sigma_j}$ separates the finitely many distinct cosets
in (E.177) in a finite quotient of $G$.  Intersect these
finitely many normal subgroups for all $i,j$.  The resulting
$N$ satisfies $N\cap C_{\sigma_j}\subseteq M_j$ and keeps
all the distinct cosets in (E.177) distinct modulo $N$.

Suppose two arrows have the same image in (E.176).  The coset
separation first puts them in the same summand, so write them
$c g_\rho\Delta_{\sigma_i}$ and
$c'g_\rho\Delta_{\sigma_i}$.  Their equality modulo $N$
says that for some
$d\in g_\rho\Delta_{\sigma_i}g_\rho^{-1}$,
$c'^{-1}c d^{-1}\in N\cap C_{\sigma_j}\subseteq M_j$.
The factor $M_j$ acts trivially modulo $R$, and $d$ fixes the
chosen arrow.  The two arrows are therefore equal modulo $R$.
Thus $C_N$ refines $R$, proving the claimed cofinality on all
Hom-sets at once, with all composition maps retained.

Taking the inverse limit, using the full topology and
separability already proved, gives the explicit formula
\[
 \widehat{\Hom}(i,j)=
 \coprod_\rho
 \widehat C_{\sigma_j}g_\rho/\widehat\Delta_{\sigma_i}
          \lhook\joinrel\longrightarrow
                 \widehat G/\widehat\Delta_{\sigma_i}.         \tag{E.178}
\]
Each summand embeds because its conjugate twist subgroup is
contained in $\widehat C_{\sigma_j}$; distinct summands remain
distinct by the same finite coset separators.  The map is
composition with the chosen $[1]:j\to\varnothing$.  Also the
completed endpoint group is exactly
$\widehat C_{\sigma_j}/\widehat\Delta_{\sigma_j}$.

If an equivalence acts trivially on the open group, its action
on the left $\widehat G$-set of crossings from $i$ to the open
object is given by
$N_{\widehat G}(\widehat\Delta_{\sigma_i})/
\widehat\Delta_{\sigma_i}$.  By (E.170) this is precisely the
endpoint automorphism group, acting freely and transitively
on the possible images of its chosen crossing.  Gauge by that
unique endpoint element to fix the crossing.  Formula (E.178)
then forces every arrow to be fixed.  The only possible
remaining natural self-comparison is the center of the open
group; it is trivial for $n\geq4$, and the $n=3$ category has
one object and trivial group.  The standard GT open-group
classification and coherent braid formulas in E.16 now give
its stated full-system comparison, with all the formerly
implicit completion steps supplied by (E.170)--(E.178).
\end{proof}
